\begin{ouverture}[calcul]{4}{Additionner et soustraire des relatifs}{Partie 1 · Nombres et calculs}
{Déplacements · Opposés · Sommes algébriques}
\objectif{Additionner des décimaux relatifs.}
\objectif{Transformer une soustraction en addition.}
\objectif{Interpréter une différence et une variation.}
\prerequis{Ordre, opposés et valeur absolue.}
\end{ouverture}
\pagestyle{otMain}\essentiel\fontsize{11}{15}\selectfont
\phantomsection\label{cours-4-1}
\essSection{1}{Construire l’addition}
\begin{propriete}Pour additionner deux nombres de même signe, additionner les distances à zéro et garder le signe commun. S’ils sont de signes contraires, soustraire la plus petite distance à la plus grande et garder le signe du nombre le plus éloigné de zéro. Deux opposés ont pour somme 0.\end{propriete}
\label{fig-cours-4-1-cours-addition}\figureGeo{cours-addition}

\begin{exemple}$(-4)+(-3)=-7$ ; $(-8)+5=-3$ ; $(-2)+7=5$. Sur une droite, ajouter un positif déplace vers la droite ; ajouter un négatif vers la gauche.\end{exemple}
\begin{savoirfaire}[Regrouper pour simplifier]On peut changer l’ordre des termes et les regrouper dans une addition : $-7+4+7-2=(-7+7)+(4-2)=2$.\end{savoirfaire}
\clearpage
\phantomsection\label{cours-4-2}
\essSection{2}{Soustraire et écrire clairement}
\begin{propriete}Soustraire un nombre revient à ajouter son opposé : $a-b=a+(-b)$. Ainsi $3-(-5)=3+5=8$, tandis que $3-5=3+(-5)=-2$.\end{propriete}
\begin{savoirfaire}[Simplifier sans perdre un signe]$(-4)+(+6)-(-2)=-4+6+2=4$. On supprime les parenthèses en tenant compte de l’opération qui les précède. Dans $8-(3-5)$, calculer d’abord la parenthèse ou changer les deux signes : $8-3+5=10$.\end{savoirfaire}
\begin{exemple}Une température passe de $-6$ à 2 degrés. Sa variation est valeur finale moins valeur initiale : $2-(-6)=8$ degrés. La température finale et la hausse sont deux quantités différentes.\end{exemple}
\begin{attention}Une différence n’est pas une distance automatiquement positive : $-2-5=-7$. La distance entre les points d’abscisses $-2$ et 5 vaut 7. Le contexte indique ce qu’il faut calculer.\end{attention}

\clearpage\rubrique{Découverte}
\begin{blocLivre}{ex-4-1}\begin{exercice}[1]{Même signe positif}
Calculer $3+8$.
\end{exercice}\end{blocLivre}
\begin{blocLivre}{ex-4-2}\begin{exercice}[1]{Même signe négatif}
Calculer $(-4)+(-9)$.
\end{exercice}\end{blocLivre}
\begin{blocLivre}{ex-4-3}\begin{exercice}[1]{Signes contraires}
Calculer $-12+7$.
\end{exercice}\end{blocLivre}
\begin{blocLivre}{ex-4-4}\begin{exercice}[1]{Soustraire un négatif}
Calculer $6-(-8)$.
\end{exercice}\end{blocLivre}
\begin{blocLivre}{ex-4-5}\begin{exercice}[1]{Deux opposés}
Calculer $-3{,}7+3{,}7$.
\end{exercice}\end{blocLivre}

\rubrique{Entraînement}
\begin{blocLivre}{ex-4-6}\begin{exercice}[2]{Deux décimaux}
Calculer $-2{,}8+4{,}15$ et $-1{,}6-2{,}9$.
\end{exercice}\end{blocLivre}
\begin{blocLivre}{ex-4-7}\begin{exercice}[2]{Réorganiser une somme}
Calculer $-17+6+17-4-6$.
\end{exercice}\end{blocLivre}
\begin{blocLivre}{ex-4-8}\begin{exercice}[2]{Parenthèses utiles}
Calculer $5-(-2+7)$ puis $5-(-2)+7$.
\end{exercice}\end{blocLivre}
\begin{blocLivre}{ex-4-9}\begin{exercice}[2]{Écrire plus simplement}
Simplifier puis calculer $(-8)-(-3)+(+6)-(+4)$.
\end{exercice}\end{blocLivre}
\begin{blocLivre}{ex-4-10}\begin{exercice}[2]{Une erreur de signe}
Inès écrit $-5-(-9)=-14$. Corriger et vérifier.
\end{exercice}\end{blocLivre}

\rubrique{Problèmes}
\begin{blocLivre}{ex-4-11}\begin{exercice}[3]{Températures}
Il fait $-4{,}5\,^\circ$C le matin. La température monte de 7 degrés puis baisse de 3,2 degrés. Quelle est la température finale ?
\end{exercice}\end{blocLivre}
\begin{blocLivre}{ex-4-12}\begin{exercice}[3]{Un ascenseur}
Un ascenseur part de l’étage $-3$, monte de huit niveaux, descend de six puis monte de deux. Où arrive-t-il ? Quelle est sa variation totale ?
\end{exercice}\end{blocLivre}
\begin{blocLivre}{ex-4-13}\begin{exercice}[3]{Distance sur une droite}
A et B ont pour abscisses $-6{,}5$ et $2{,}25$. Calculer la distance AB.
\label{fig-ex-4-13}\figureGeo{add-13}
\end{exercice}\end{blocLivre}
\begin{blocLivre}{ex-4-14}\begin{exercice}[3]{Retrouver une variation}
Un compte fictif passe de $-18{,}40$ euros à 7,25 euros après une seule opération. Quelle opération a eu lieu ?
\end{exercice}\end{blocLivre}
\begin{blocLivre}{ex-4-15}\begin{exercice}[3]{Un bilan à reconstituer}
Un jeu commence avec un score inconnu. On perd 12 points, on gagne 7 puis on perd 9. Le score final est $-6$. Retrouver le score initial.
\end{exercice}\end{blocLivre}

\clearpage
